\mytitle{\thedaynumber/\death}
\hypertarget{\the\value{daynumber}}

\marginnote[1mm]{\textit{The right information in the right place just changes your life for the better.
And if it's so valuable, people will have a desire to share the source with their circle.}}

{\HUGE \theexcluded\ Drafts} \stepcounter{excluded}

Degrees of freedom
When tech shows up, it offers a shortcut and convenience.

You can use Google Maps to direct you somewhere without paying much attention to the surroundings.

You can use Claude to write your marketing copy and get a better-than-mediocre result the first try.

You can look for a gift on Amazon, pick the first match, and be pretty sure it’ll do the job.

Tech adoption often focuses on making things easier, simpler, and pre-decided.

And yet… we can also decide to use tech to do more work, insert more humanity, and amplify flexibility. We don’t try
to get our time back, we try to figure out how to leverage the time we’ve got.

When a film director uses AI to create storyboards, it’s a chance to generate multiple approaches to a scene, not just one. When we sit with all the data Google Maps offers us for a trip, we might plan a less direct route, with more stops and detours, simply because we now know what our options are. And once we know what the mediocre and average marketing copy looks like, we put in the time (and take the risks) to go to edges we never would have had the resources to explore in the old days.

The best tech gives us a chance to work harder on the parts that matter to our customers and to us.

\textit{Here’s the simple fork in the road:}

Professionals and organizations that use AI to save time, cut costs, and lay people off are taking a lazy road to failure and irrelevance.

Those who use it to do harder, braver, and more powerful work, who figure out how to create more value and charge more for it, and who end up hiring more people to do so, will be defining our future.


